\documentclass[11pt]{article}
\usepackage[margin=1in]{geometry}
\usepackage{enumitem}
\usepackage{hyperref}
\hypersetup{colorlinks=true, urlcolor=blue, linkcolor=blue}
\usepackage{parskip}

\begin{document}

\begin{center}
  {\Large\bfseries Exercise 0: Your First Scrape}\\[2pt]
  {\normalsize SISMID 2026 \textbullet\ Day 1, 9:45 \textbullet\ Agent Coding Introduction}
\end{center}

\medskip

Before we model any disease signal, we have to \emph{get} one. In this warm-up you set
up an identical coding environment, stand up an AI coding agent, and use it (or a
pre-filled notebook) to scrape your first novel data stream: \textbf{Google Dengue
Trends for Mexico}. This is the same search signal you will model against reported
dengue cases in the very next session, so the data you produce here carries straight
into the 11:00 exercise. You do it in two parts: a small \emph{live} pull of a handful
of terms, to learn the flow and see present-day searches, and then the \emph{historical}
2004--2011 window that feeds the model.

We start with Google Trends on purpose. It is the backbone of this whole course, and it
is also \emph{awkward}: there is no official API and it actively rate-limits scrapers,
so a live pull often fails and returns different numbers each time. That awkwardness is
a lesson, not a bug, and it sets up this afternoon's session on preprocessing Trends.

You may work in \textbf{Python or R}. Two lanes lead to the same place, and you may
switch at any time:
\begin{itemize}[nosep]
  \item \textbf{Lane A (agent).} Drive a coding agent (Codex, Claude Code, or Antigravity
        CLI) with the prompts in \texttt{notebooks/01\_data\_scraping\_intro.ipynb}.
  \item \textbf{Lane B (notebook).} Run the pre-filled
        \texttt{notebooks/01\_data\_scraping\_intro\_soln.ipynb} cell by cell.
\end{itemize}

\textbf{Materials.} Course repository:
\url{https://github.com/Shihao-Yang/sismid2026}. Cached fallback data for this session:
\texttt{data/google\_trends\_dengue\_mx\_cached.csv} (used automatically if a live pull
is rate-limited).

\begin{enumerate}[label=(\alph*)]
  \item \textbf{Set up an identical environment.} Open the course repository and launch
        it one of three ways: (i) a \textbf{GitHub Codespace} from the browser
        (\texttt{Code} $\rightarrow$ \texttt{Codespaces} $\rightarrow$ \texttt{Create});
        (ii) \textbf{clone} the repository and work locally; or (iii) open the notebooks
        in \textbf{Google Colab}. Run the smoke test
        (\texttt{notebooks/00\_smoke\_test.ipynb}) and confirm you see the library
        versions and ``Environment OK''.

  \item \textbf{Stand up the coding agents.} In a terminal, run
        \texttt{bash scripts/setup-agents.sh} to install the Codex, Claude Code, and Antigravity
        CLIs, then authenticate using the method the instructors provide in class. Sanity
        check: ask an agent to ``list the CSV files in \texttt{data/} and print the first
        rows of each.'' (If you cannot get an agent working, use Lane B and continue; you
        will not be blocked.)

  \item \textbf{See it live (a handful of terms).} Using \texttt{pytrends}, pull a
        \emph{handful} of dengue-related terms (\emph{dengue}, \emph{sintomas de dengue},
        \emph{mosquito}) for \textbf{Mexico} (\texttt{geo=MX}) over the past five years
        (\texttt{today 5-y}). Confirm the last data point lands on the \textbf{current
        week}: this is a genuinely live pull. Keep it to a handful of terms, which is what
        stays under Google's rate limit, and plot the three series.

  \item \textbf{Historical window for 11:00.} Now pull the four terms (\emph{dengue},
        \emph{sintomas de dengue}, \emph{mosquito}, \emph{dengue sintomas}) for Mexico over
        \textbf{2004--2011}, the same period as the reported-case ground truth you model
        next session. If a live pull is rate-limited, fall back to the cached CSV. Merge on
        date and overlay all four. Which terms look like they track disease activity, and
        which look like noise? (You will let least squares decide at 11:00.)

  \item \textbf{Note the instability, then save.} Pull the single term \emph{dengue}
        \textbf{twice} over the recent window. Google Trends is built from a sample of
        searches, so pulls can differ: usually only slightly on the public 0--100 endpoint,
        but substantially on the raw Health Trends API we discuss after lunch. Note what you
        see, then save your historical four-term table as a tidy CSV (one date column, one
        column per term). This is the file you carry into 11:00.

  \item \textbf{Discuss.} Google Trends measures \emph{attention}, not infection, and it
        is not reproducible from one pull to the next. Name two ways this could mislead a
        naive model, and one thing you could do about the reproducibility problem.
        (Hint: what happens if you average many pulls? We return to this after lunch.)
\end{enumerate}

\vfill
\noindent\textit{Stretch:} change \texttt{geo} to your own country, or swap the keywords
for a disease you care about, and rerun. If you are on Lane A, do it by editing a single
sentence of your prompt.

\end{document}
